\subsection{Conjectures}

\note{I would like to produce a table of the minimum number of additional edges from attempting to subset $\Grid_{k \times n}$ from $\CGrid$. The cases $k=2,3$ already appear in the previous sections, but $k=4,5$ should be within reach in a few weeks. This is because of the following:}

\begin{conjecture}
\label{conjecture:min-edges}
\note{This needs to be formalized:} Suppose we have a construction to subset a $k \times n$ subgraph $H$ of some $\CGrid$ with $a$ additional edges compared to $\Grid_{k \times n}$. Suppose further than this construction produces $b_k$ additional edges for every column added, and this $b$ remains constant regardless of how many additional columns are added. Then a subgraph $H'$ of any $\CGrid$ with additional edges compared to $\Grid_{k+1 \times n}$ must have at least 
\begin{align}
    b_{k+1} \geq \lceil b_k \frac{k+1}{k} \rceil
\end{align}
additional edges for every column added.
\begin{proof}
By contradiction and induction. The base case is $3 \times n$ where $b_3 = 1$, and $4 \times n$ where if $b_4 = 1$, then at least two rows have no additional edge. Deleting the one row connected to this additional edge \note{and performing some transformations afterwards, which need to be formalized} leaves a $3 \times n$ grid with fewer than $b_3$ additional edges per column. Contradiction. \note{It is not quite clear whether this generalizes and hence is left as a conjecture.}
\end{proof}
\end{conjecture}

\begin{conjecture}
Using the same notation as Conjecture \ref{conjecture:min-edges}, whatever construction $H$ is minimal for $k$ rows has $H$ be a bipartite graph. In particular, it is of the form of two antidiagonal regions of the comparability grid. \note{Since the bipartite graphs with $n$ vertices are counted by $C_n$, the $n$th Catalan number, which is large but smaller than $n!$, significant early termination is still necessary to efficiently find the minimum number of additional edges for any $k$}.
\end{conjecture}

It is a fact that the number of additional edges must increase as $k$ increases. In that sense, the $3 \times n$ case is the easiest case to check universality through algorithmic approaches for. 

\note{Here I want to provide a table of the results of the size of the groups generated using the width-3-gadget, and if there is enough compute, the width-4-gadget.}

Another algorithmic direction is to include local complementation into the algorithmic search. \note{However this is computationally quite expensive, because it means we cannot use the tuple subset representation of Fact \ref{fact:tuple-construction}. Furthermore, (local) graph isomorphism is now necessary to prevent checking the same graph multiple times, which further adds to the complexity of the algorithm}.